\chapter{La Regola dell'Ottava: versione completa}
\label{chap:regola_ottava_completa}

Dopo aver esplorato la teoria generale delle settime nel Capitolo \ref{chap:sette_specie}, presentiamo la \textbf{versione completa} della Regola dell'Ottava, che integra tutte le settime diatoniche nell'armonizzazione della scala al basso.

\section{Le tre versioni a confronto}
\label{sec:tre_versioni_confronto}

\begin{definizione}[Le tre versioni della Regola dell'Ottava]
	\label{def:tre_versioni_complete}
	
	\begin{enumerate}
		\item \textbf{Versione preliminare (Capitolo \ref{chap:funzioni_armoniche}):}
		\begin{itemize}
			\item Utilizza solo \textbf{triadi} in stato fondamentale e rivolti
			\item Armonizzazione "sicura" e consonante
			\item Ideale per l'apprendimento iniziale della condotta delle voci
		\end{itemize}
		
		\item \textbf{Versione intermedia (Capitolo \ref{chap:settima_dominante}):}
		\begin{itemize}
			\item Aggiunge la \textbf{settima di dominante (\sigla{V}{7}{})} e i suoi tre rivolti
			\item Introduce la dissonanza controllata
			\item Maggiore tensione e direzionalità armonica
		\end{itemize}
		
		\item \textbf{Versione completa (questo capitolo):}
		\begin{itemize}
			\item Include \textbf{tutte le settime diatoniche}: \sigla{II}{7}{}, \sigla{IV}{7}{}, \sigla{V}{7}{}, \sigla{VII}{7}{}
			\item Massima ricchezza armonica e varietà di colori
			\item Riflette la pratica compositiva del tardo Barocco e primo Classicismo
		\end{itemize}
	\end{enumerate}
\end{definizione}

\begin{matematica}[Crescita combinatoria delle possibilità armoniche]
	Analizziamo quantitativamente come cresce la complessità passando da una versione all'altra:
	
	\textbf{Versione preliminare:}
	\begin{itemize}
		\item 7 triadi diatoniche
		\item 3 stati per ogni triade (fondamentale + 2 rivolti)
		\item $7 \times 3 = 21$ possibili configurazioni di base
	\end{itemize}
	
	\textbf{Versione intermedia:}
	\begin{itemize}
		\item 21 configurazioni precedenti
		\item +1 accordo di settima (\sigla{V}{7}{}) con 4 stati (fondamentale + 3 rivolti)
		\item $21 + 4 = 25$ configurazioni disponibili
	\end{itemize}
	
	\textbf{Versione completa:}
	\begin{itemize}
		\item Configurazioni precedenti
		\item +6 nuove settime diatoniche (\sigla{I}{7}{}, \sigla{II}{7}{}, \sigla{III}{7}{}, \sigla{IV}{7}{}, \sigla{VI}{7}{}, \sigla{VII}{7}{}) × 4 stati ciascuna
		\item Totale teorico: $21 + (7 \times 4) = 49$ configurazioni
	\end{itemize}
	
	La versione completa offre più del doppio delle possibilità armoniche rispetto alla versione preliminare.
\end{matematica}

\section{La Regola dell'Ottava completa nel modo maggiore}
\label{sec:regola_completa_maggiore}

Presentiamo ora lo schema completo per il modo maggiore, integrando tutte le varianti con settime.

\begin{regola}[Regola dell'Ottava Completa - Modo Maggiore]
	\label{reg:regola_ottava_completa_maggiore}
	
	\textbf{Scala ascendente (esempio in Do maggiore):}
	
	\begin{tabular}{|c|c|c|c|}
		\hline
		\textbf{Grado} & \textbf{Basso} & \textbf{Numerica} & \textbf{Funzione armonica} \\
		\hline
		I & Do & 5 & I (Tonica) \\
		II & Re & 6 \textit{oppure} 4/3 & \sigla{VII}{6}{} \textit{oppure} \sigla{V}{4}{3} \\
		III & Mi & 6 & \sigla{I}{6}{} (Tonica) \\
		IV & Fa & 6 \textit{oppure} 6/5 & \sigla{II}{6}{} \textit{oppure} \sigla{II}{6}{5} (Sopratonica 7) \\
		V & Sol & 5 \textit{oppure} 7 & V \textit{oppure} \sigla{V}{7}{} (Dominante) \\
		VI & La & 6 & \sigla{IV}{6}{} (Sottodominante) \\
		VII & Si & 6 \textit{oppure} 6/5 & \sigla{V}{6}{} \textit{oppure} \sigla{V}{6}{5} (Dominante 7) \\
		VIII & Do & 5 & I (Tonica) \\
		\hline
	\end{tabular}
	
	\vspace{0.5cm}
	
	\textbf{Scala discendente (esempio in Do maggiore):}
	
	\begin{tabular}{|c|c|c|c|}
		\hline
		\textbf{Grado} & \textbf{Basso} & \textbf{Numerica} & \textbf{Funzione armonica} \\
		\hline
		VIII & Do & 5 & I (Tonica) \\
		VII & Si & 6 \textit{oppure} 6/5 & \sigla{V}{6}{} \textit{oppure} \sigla{V}{6}{5} (Dominante 7) \\
		VI & La & 4/3\# & \sigla{V}{4}{3}\# di V (dominante secondaria, cap.~\ref{chap:tonicizzazione}; la vera sesta eccedente richiede La$\flat$, \ref{sec:sesta_eccedente}) \\
		V & Sol & 5 \textit{oppure} 7 & V \textit{oppure} \sigla{V}{7}{} (Dominante) \\
		IV & Fa & 2 & \sigla{V}{2}{} (Dominante terza inversione) \\
		III & Mi & 6 & \sigla{I}{6}{} (Tonica) \\
		II & Re & 4/3 & \sigla{V}{4}{3} (Dominante seconda inversione) \\
		I & Do & 5 & I (Tonica) \\
		\hline
	\end{tabular}
\end{regola}

\begin{esempio}[Regola dell'Ottava Completa in Do Maggiore - scala discendente con tutte le settime]
	\label{ex:regola_completa_disc_magg}
	\musica{es_cap14_01}
	Nota: Sul VI grado discendente (La), la cifratura 4/3\# indica la dominante secondaria \sigla{V}{4}{3} di V (Re-Fa\diesis{}-La-Do con il La al basso). Non è una sesta eccedente: quella in maggiore si forma sul La$\flat$ (VI abbassato), vedi il capitolo \ref{chap:seste_eccedenti}.
\end{esempio}

\section{La Regola dell'Ottava completa nel modo minore}
\label{sec:regola_completa_minore}

Nel modo minore, la versione completa presenta alcune particolarità, legate alla scala minore armonica.

\begin{regola}[Regola dell'Ottava Completa - Modo Minore]
	\label{reg:regola_ottava_completa_minore}
	
	\textbf{Scala ascendente (esempio in La minore armonico):}
	
	\begin{tabular}{|c|c|c|c|}
		\hline
		\textbf{Grado} & \textbf{Basso} & \textbf{Numerica} & \textbf{Funzione armonica} \\
		\hline
		I & La & 5 & i (Tonica) \\
		II & Si & 4/3 & \sigla{V}{4}{3} (Dominante) \\
		III & Do & 6 & \sigla{i}{6}{} (Tonica) \\
		IV & Re & 6 \textit{oppure} 6/5 & \sigla{ii°}{6}{} \textit{oppure} \sigla{ii°}{6}{5} (Sopratonica 7) \\
		V & Mi & 5 \textit{oppure} 7 & V \textit{oppure} \sigla{V}{7}{} (Dominante) \\
		VI & Fa & 6 & \sigla{iv}{6}{} (Sottodominante) \\
		VII & \Soldiesis  & 6/5 & \sigla{V}{6}{5} (Dominante 7) \\
		VIII & La & 5 & i (Tonica) \\
		\hline
	\end{tabular}
	
	\vspace{0.5cm}
	
	\textbf{Scala discendente (esempio in La minore):}
	
	\begin{tabular}{|c|c|c|c|}
		\hline
		\textbf{Grado} & \textbf{Basso} & \textbf{Numerica} & \textbf{Funzione armonica} \\
		\hline
		VIII & La & 5 & i (Tonica) \\
		VII & \Soldiesis  & 6 & \sigla{V}{6}{} (Dominante) \\
		VI & Fa & \#6 & iv\#6 (Sesta Eccedente, \ref{sec:sesta_eccedente}) \\
		V & Mi & 5 \textit{oppure} 7 & V \textit{oppure} \sigla{V}{7}{} (Dominante) \\
		IV & Re & 2 & \sigla{V}{2}{} (Dominante terza inversione) \\
		III & Do & 6 & \sigla{i}{6}{} (Tonica) \\
		II & Si & 4/3 & \sigla{V}{4}{3} (Dominante seconda inversione) \\
		I & La & 5 & i (Tonica) \\
		\hline
	\end{tabular}
\end{regola}

\begin{esempio}[Regola dell'Ottava Completa in La minore - scala ascendente]
	\label{ex:regola_completa_asc_min}
	\musica{es_cap14_02}
\end{esempio}

\begin{esempio}[Regola dell'Ottava Completa in La minore - scala discendente]
	\label{ex:regola_completa_disc_min}
	\musica{es_cap14_03}
	Nota: Sul VI grado discendente (Fa), la cifratura \#6 (sesta eccedente) indica la caratteristica alterazione ascendente della sesta (\Rediesis ) che risolve espandendosi verso l'ottava (Mi).
\end{esempio}

\section{La sesta eccedente nel contesto della Regola dell'Ottava}
\label{sec:sesta_eccedente}

Nelle tabelle della Regola dell'Ottava completa, sul VI grado discendente appare spesso la cifratura $\sharp 6$ o $4/3\sharp$ con l'etichetta "sesta eccedente".

\begin{definizione}[Accordo di sesta eccedente]
    \label{def:sesta_eccedente}
    È un accordo cromatico caratterizzato dall'intervallo di sesta eccedente tra il basso e una voce superiore. Si forma tipicamente tra il VI grado abbassato al basso e il IV grado alzato cromaticamente in una voce superiore.
    La sua funzione principale è preparare la dominante con una spinta cromatica convergente verso la nota della dominante stessa.
\end{definizione}

\subsection{Le tre specie di sesta eccedente}

\begin{itemize}
    \item \textbf{Sesta italiana (It+6):} Formata da VI$\flat$ (basso), tonica e IV$\sharp$ (sesta eccedente). Solo tre note.
    \item \textbf{Sesta francese (Fr+6):} Aggiunge il II grado alla sesta italiana. In Do m: \Labemolle -Do-Re-\Fadiesis .
    \item \textbf{Sesta tedesca (Ger+6):} Aggiunge il III$\flat$ alla sesta italiana. In Do m: \Labemolle -Do-\Mibemolle -\Fadiesis . Produce quinte parallele se risolve direttamente su V; si usa spesso il \sigla{I}{6}{4} come ponte (\ref{chap:quarta_sesta}).
\end{itemize}

Le tre specie, con struttura, condotta delle voci, tabelle in tutte le tonalità ed errori da evitare, sono trattate nel capitolo \ref{chap:seste_eccedenti}. Ecco la costruzione in La minore (con il Fa al basso) e in Do maggiore (con il La$\flat$ al basso):
\musica{es_cap22_01}
\musica{es_cap22_02}

\begin{regola}[Risoluzione della sesta eccedente]
    \label{reg:risoluzione_sesta_eccedente}
    \begin{itemize}
        \item Il basso (VI$\flat$) scende di semitono alla dominante.
        \item La voce con la sesta eccedente (IV$\sharp$) sale di semitono alla dominante.
        \item Le altre voci si muovono per grado congiunto verso il V o verso il \sigla{I}{6}{4} preparatorio.
    \end{itemize}
\end{regola}

La sesta eccedente appare naturalmente nella Regola dell'Ottava sul VI grado discendente, e rafforza la spinta verso la dominante (V grado).

\section{Usi pratici e stilistici}
\label{sec:usi_pratici_completi}

La versione completa della Regola dell'Ottava offre \textbf{possibilità} tra cui il compositore sceglie in base al contesto musicale.

\begin{regola}[Criteri di scelta tra le versioni]
	\label{reg:criteri_scelta_versioni}
	
	\begin{itemize}
		\item \textbf{Periodo storico:}
		\begin{itemize}
			\item Primo Barocco: preferire la versione preliminare
			\item Tardo Barocco (Bach, Händel): versione intermedia o completa
			\item Classicismo (Mozart, Haydn): versione completa con libertà di scelta
		\end{itemize}
		
		\item \textbf{Contesto espressivo:}
		\begin{itemize}
			\item Passaggi leggeri, danzanti: versione preliminare
			\item Momenti drammatici, intensi: versione intermedia o completa
			\item Sequenze modulanti: versione completa (le settime facilitano la modulazione)
		\end{itemize}
		
		\item \textbf{Velocità del movimento del basso:}
		\begin{itemize}
			\item Basso veloce (semicrome, terzine): preferire triadi semplici
			\item Basso moderato (crome, semiminime): versione intermedia
			\item Basso lento (minime, semibrevi): versione completa con massima ricchezza
		\end{itemize}
		
		\item \textbf{Texture:}
		\begin{itemize}
			\item Contrappunto denso: versione preliminare (evitare troppe dissonanze)
			\item Accompagnamento omofono: versione completa (maggiore libertà armonica)
		\end{itemize}
	\end{itemize}
\end{regola}

\begin{fisica}[Perché le settime "colorano" l'armonia?]
	Dal punto di vista acustico, l'aggiunta di settime agli accordi produce effetti percettivi precisi:
	
	\begin{enumerate}
		\item \textbf{Rugosità timbrica:} Gli intervalli di settima (maggiore o minore) creano battimenti tra armonici parziali, producendo una texture sonora più "densa" e complessa
		
		\item \textbf{Fusione incompleta:} A differenza delle triadi, dove le note tendono a "fondersi" in un'unica sensazione sonora, le settime mantengono una certa separazione percettiva, creando un senso di "movimento interno"
		
		\item \textbf{Tensione direzionale:} Le dissonanze di settima attivano aspettative percettive di risoluzione, guidando l'orecchio verso specifiche destinazioni armoniche
	\end{enumerate}
\end{fisica}

\section{Esercizi di armonizzazione completa}
\label{sec:esercizi_armonizzazione}

Per padroneggiare la versione completa della Regola dell'Ottava, serve esercitarsi con bassi dati che richiedono scelte armoniche consapevoli.

\begin{esempio}[Basso dato con armonizzazione mista]
	\label{ex:basso_dato_misto}
	
	Armonizzare il seguente basso scegliendo strategicamente tra le tre versioni della Regola dell'Ottava:
	
	\musica{es_cap14_07}
	
	\textbf{Soluzione suggerita:}
	\begin{itemize}
		\item Battute 1-2 (scala ascendente e discendente): usare versione completa per massima varietà
		\item Battuta 3 (salto Do-Sol): usare stati fondamentali (I-V)
		\item Battuta 4 (cadenza finale): ritorno alla semplicità della versione preliminare
	\end{itemize}
\end{esempio}

% ============================================================
% FINE NUOVO CAPITOLO REGOLA DELL'OTTAVA COMPLETA
% ============================================================


% ============================================================
% TABELLA RIASSUNTIVA DELLA REGOLA DELL'OTTAVA COMPLETA
% Da inserire come appendice o tabella di riferimento rapido
% ============================================================

\begin{table}[h]
	\centering
	\caption{Regola dell'Ottava Completa - Tabella di riferimento}
	\label{tab:regola_ottava_completa_riepilogo}
	\small
	\begin{tabular}{|c|c|c|c|c|c|}
		\hline
		\textbf{Dir.} & \textbf{Grado} & \textbf{Basso} & \textbf{MAGGIORE} & \textbf{MAGGIORE} & \textbf{MINORE} \\
		& \textbf{Scala} & \textbf{(Es. Do/La)} & \textbf{Numerica} & \textbf{Funzione} & \textbf{Numerica} \\
		\hline
		Asc & I & Do / La & 5 & I & 5 \\
		Asc & II & Re / Si & 6 (o 4/3) & \sigla{VII}{6}{} (o \sigla{V}{4}{3}) & 4/3 \\
		Asc & III & Mi / Do & 6 & \sigla{I}{6}{} & 6 \\
		Asc & IV & Fa / Re & 6 (o 6/5) & \sigla{II}{6}{} (o \sigla{II}{6}{5}) & 6 (o 6/5) \\
		Asc & V & Sol / Mi & 5 (o 7) & V (o \sigla{V}{7}{}) & 7 \\
		Asc & VI & La / Fa & 6 & \sigla{IV}{6}{} & 6 \\
		Asc & VII & Si / \Soldiesis  & 6 (o 6/5) & \sigla{V}{6}{} (o \sigla{V}{6}{5}) & 6/5 \\
		Asc & VIII & Do / La & 5 & I & 5 \\
		\hline
		Disc & VIII & Do / La & 5 & I & 5 \\
		Disc & VII & Si / Sol & 6 (o 6/5) & \sigla{V}{6}{} (o \sigla{V}{6}{5}) & 6 \\
		Disc & VI & La / Fa & 4/3\# & \sigla{II}{4}{3}\# (V del V) & \#6 \\
		Disc & V & Sol / Mi & 5 (o 7) & V (o \sigla{V}{7}{}) & 5 (o 7) \\
		Disc & IV & Fa / Re & 2 & \sigla{V}{2}{} & 2 \\
		Disc & III & Mi / Do & 6 & \sigla{I}{6}{} & 6 \\
		Disc & II & Re / Si & 4/3 & \sigla{V}{4}{3} & 4/3 \\
		Disc & I & Do / La & 5 & I & 5 \\
		\hline
	\end{tabular}
\end{table}
